\chapter{Planetary, satellite and comet applications}\label{ch:planetary}
\section{A common kinetic engine, different physical targets}
Planetary applications choose their own body shape, rotation, gravity,
illumination, surface inventories, plasma/radiation fields and observation
operators. A species list or successful mesh initialization is not a complete
physical model. Record the target, epoch, frame, source law and loss/network
scope with every run. Use ephemerides and orientation consistently in
injection, acceleration, surface interception and synthetic observations.

\section{The Moon: exosphere, terrain and observation geometry}
\path{srcMoon} combines spherical or terrain-based surface geometry with
target-specific exosphere processes and diagnostics. Its generic runtime
block is selected with \code{-input FILE}; omission preserves the historical
analytic-sphere regression. The terrain path requires declared LOLA products,
labels, resolution and output paths. Do not replace unavailable topography
with a sphere without recording a different physical case.

Spatial targets distinguish the near-surface, far-field, terminator and polar
regions. Quick, regression and full science receipts are different numerical
campaigns. A completed coarse wiring run must retain its M1
SKIPPED/NOT QUALIFIED designation when convergence criteria were not exercised.
The lunar unit/integration runners cover geometry, source/thermal state,
environment timing and bounded products; inspect their actual per-case status.
Na tail and limb observations additionally require instrument geometry,
velocity-dependent radiation, line-of-sight weighting and the declared
exposure average.
\source{srcMoon/README.md; srcMoon/test/README.md; srcMoon/examples/README.md}

\section{Enceladus: sources, parent loss and Cassini comparisons}
\path{srcEnceladus} separates a fitted-source exporter from the forward kinetic
application. The currently implemented distribution identifier is
\code{tenishev2014_hybrid_8}: eight point jets and seven stripe components.
Stripe vertices are projected vertically onto the declared 252100-m reference
sphere. The separate source file states geometry, kinematics, frame and
whether strengths are absolute rates or a normalized shape.

Equal-strength E3/E5 receipts exist to start and exercise the forward code.
They are not a recovered individual-source fit. The E3/E5 mass-44
normalizations are apparent CO2 and carry a documented impact-fragmentation
bracket; they must not be reported as unqualified physical CO2 production.
The schema-2 run transports neutral H2O/CO2 with selected H2O photolytic
parent loss only, no daughter particles and no plasma.

Observation requests select INMS density/count products and supported UVIS
occultations. Density, detector count rate and occultation transmission are
different operators. Only compare like quantities with compatible response,
units, frame and epoch. A trajectory outside the mesh has no local density;
do not interpolate across that gap or label it zero. Resolve closest approach
with actual flyby cells, and plot sampling uncertainty and data availability.
\source{srcEnceladus/SOURCE_DISTRIBUTION_FILE.md; srcEnceladus/test/README.md}

\section{Comets: radial and shape-based gas/dust models}
\path{srcHartley2} retains its historical name but contains a generic comet
layer with a DIM1 spherical radial backend, legacy DIM3 shape application,
configurable chemistry/collisions, internal energy, cooling and brightness.
DIM2 axisymmetric execution is reserved rather than implemented. I3 uses
\path{input/i3.conf} for compiled capabilities and \path{input/i3.in} plus
\code{-case N} for physical cases. Halley and Hartley decks use their
documented workflows; do not assume every alias stages the shared parser
identically.

The radial backend retains three-dimensional velocity/trajectory geometry
while exploiting spherical symmetry for storage and sampling. It cannot
represent an arbitrary asymmetric jet or shadowed nonspherical coma.
For those cases use the appropriate DIM3 adapter with a qualified surface,
thermal map and gravity model. Nucleus production and escaping/net production
can differ because of returning particles; state which quantity is fitted.

The separate collisionless reference supplies a far-coma column correction
for selected I3 post-processing. It is not an extension of the DSMC domain.
Its columns, normalization, assumptions and missing radiation-pressure effects
are documented independently. Comparison tables imported from COMET1D are
historical evidence, not new runs of this checkout.
\source{srcHartley2/README.md; srcHartley2/COLLISIONLESS_REFERENCE.md; srcHartley2/i3/COMPARISON.md}

\section{Planetary energetic particles}
\path{srcEarth} includes backward-characteristic cutoff/access and
instantaneous spectrum products using declared static/quasi-static field
snapshots. A backward access calculation and a forward Monte Carlo density
calculation are different estimators. Field, spectrum, detector response and
observer epoch must form one accepted configuration with consistent units.
Classify unresolved/trapped/escaped characteristics explicitly; an unfinished
trajectory must not silently count as excluded flux.

The coupled Step-10/11 contracts produce content-identified snapshots and
product manifests. Their portable analytic/replay tests do not qualify
all external-data or live SWMF release matrices. Magnetic-only instantaneous
results do not imply electric acceleration or long-duration validity across
changing fields. Other planetary application directories require their own
input/data and acceptance review before reuse.
\source{srcEarth/README.md}
